\documentclass[11pt,a4paper]{article}
\usepackage{fontspec}
\setmainfont{DejaVu Sans}[Scale=0.90]
\usepackage{amsmath,amssymb}
\usepackage[margin=2.2cm]{geometry}
\usepackage{booktabs,array,tabularx}
\usepackage{graphicx}
\graphicspath{{../figures/}}
\usepackage{enumitem}
\usepackage{titlesec}
\usepackage{parskip}
\titlelabel{\thetitle.\quad}
\titleformat{\section}{\normalfont\large\bfseries}{\thesection.}{0.6em}{}
\titleformat{\subsection}{\normalfont\normalsize\bfseries}{\thesubsection}{0.6em}{}
\titlespacing*{\section}{0pt}{15pt}{5pt}
\titlespacing*{\subsection}{0pt}{11pt}{4pt}
\setlist[itemize]{leftmargin=1.4em,itemsep=1pt,topsep=3pt}
% Preserve fonts/margins; permit a final line-breaking pass for long prose.
\setlength{\emergencystretch}{2em}
\begin{document}
\begin{center}
{\Large\bfseries Validation d'un opérateur par événements}\\[4pt]
{\normalsize Collisions de Bose, conservation et parité}\\[2pt]
{\small Partie I --- 24 septembre 2026}
\end{center}
\section{Question et statut}
La note 18 identifie un obstacle : une diagonale de durées de vie ou un
opérateur équivalent pour la conductivité ne suffit pas à certifier la
dynamique des populations. Cette note construit une référence finie
permettant de vérifier un futur import d'événements.
Les énergies et les taux des exemples sont synthétiques. Aucun taux
microscopique de l'AlN n'est déduit de ces exemples.

\section{Un événement réversible compté une fois}
On considère des modes distincts, des fréquences strictement positives,
une température positive et une grille complète à poids commun.
Pour $i\leftrightarrow j+k$, on pose
\[
 x_\mu=\frac{\epsilon_\mu}{k_BT},\quad x_i=x_j+x_k,\quad
 N_\mu=\frac{1}{e^{x_\mu}-1},\quad
 d_\mu=\sqrt{N_\mu(1+N_\mu)}.
\]
Le vecteur stoechiométrique vaut $s=-e_i+e_j+e_k$ et
\[
 \dot n=sF(n),\qquad
 F(n)=\Gamma\{n_i(1+n_j)(1+n_k)-(1+n_i)n_jn_k\},\quad \Gamma\ge0.
\]
La réaction inverse est déjà incluse dans cette différence. Un partenaire
obtenu par inversion des vecteurs d'onde est un autre événement, sauf
si le triplet non ordonné se retrouve identique. Les permutations des deux
filles ne doivent pas multiplier involontairement le taux.
La convention de $\Gamma$ inclut les facteurs de comptage et d'unités :
aucune conversion depuis une largeur de raie n'est supposée ici.

\section{Linéarisation et invariants}
Le bilan détaillé donne
\[
 A=\Gamma N_i(1+N_j)(1+N_k)
   =\Gamma(1+N_i)N_jN_k .
\]
Avec $D=\operatorname{diag}(d_\mu)$ et $y=D^{-1}\delta n$,
\[
 \delta F=-A\,s^TD^{-2}\delta n,\qquad
 \partial_t y=-Cy,\qquad
 C=A(D^{-1}s)(D^{-1}s)^T.
\]
La somme des événements conserve la positivité et la symétrie.
Pour toute quantité additive $h$ vérifiant $s^Th=0$ événement par événement,
\[
 C(Dh)=0,\qquad y^TCy=A|(D^{-1}s)^Ty|^2\ge0.
\]
L'énergie est conservée par résonance. L'impulsion cristalline est
conservée pour un événement normal, mais sa variation peut être un
vecteur réciproque non nul pour un événement Umklapp.
Une paire de vecteurs réciproques opposés ne fait pas disparaître
la dissipation quadratique de l'impulsion.

La Hessienne de l'entropie relative de Bose à l'équilibre est
$D^{-2}$ : la même transformation est obtenue par une deuxième voie.
Le rang est celui des vecteurs d'événements actifs. Quelques événements
laissent généralement des invariants accidentels ; on ne peut en déduire
un temps de relaxation global ni une conductivité statique finie.

\section{Pourquoi la parité doit être vérifiée}
Soit $J$ l'inversion des modes. Des partenaires de mêmes énergies
et de mêmes coefficients assurent $JC=CJ$. Un événement
$(+q,-q)\leftrightarrow0$, avec un mode optique positif en $q=0$,
peut être invariant par $J$ et agir uniquement dans le secteur pair.
Changer son coefficient peut donc préserver tout le bloc impair tout
en modifiant la dynamique complète.

Le transport spatial couple les secteurs. Pour des vitesses réciproques,
l'élimination exacte du secteur pair donne, pour $\Re z>0$,
\[
 S_{\rm impair}=z+C_{\rm impair}
  +k^2V_{\rm impair,pair}(z+C_{\rm pair})^{-1}V_{\rm pair,impair}.
\]
L'égalité des opérateurs impairs ne certifie donc pas la réponse à $k\ne0$.
Il s'agit d'un contre-exemple fini contrôlé, pas d'une mesure de l'AlN.
Son poids peut disparaître lors du passage au continuum ; aucun effet
macroscopique persistant n'est établi par cet exemple.

\section{Le contrôle négatif : événement non résonant}
Si $\Delta=x_i-x_j-x_k\ne0$, le rapport des facteurs de Bose directs et
inverses à l'équilibre vaut $e^{-\Delta}$. L'état thermique n'est plus
stationnaire pour le crochet non modifié. La production d'énergie vaut
\[
 \epsilon^T\dot n
 =-k_BT\,\Delta\,\Gamma(1+N_i)N_jN_k(e^{-\Delta}-1)\ge0 .
\]
La positivité est stricte pour $\Gamma>0$ et $\Delta\ne0$.
Les désaccords de signes opposés ne se compensent donc pas dans cette
production à l'état thermique avec des poids positifs.
Une pondération d'élargissement n'est pas une preuve de conservation.
Cette remarque ne remplace pas l'analyse d'une quadrature sur les surfaces
d'énergie du problème continu.

Une tolérance relative seule est insuffisante lorsque $x$ est très grand.
Le code contrôle également $|\Delta|$. Une somme compensée calcule le désaccord des valeurs effectivement
stockées. L'audit a trouvé que la soustraction ordinaire pouvait accepter
$(10^{16},1,10^{16})$ en annonçant zéro au lieu de $-1$.
Les deux permutations des filles sont désormais rejetées.
Cette tolérance ne représente pas une largeur spectrale physique :
un désaccord accepté mais non nul donne une approximation positive,
pas la Jacobienne exacte du crochet non stationnaire.
Aucune borne d'erreur sur une réponse inverse n'en découle.

\section{Implémentation et vérifications}
Le fichier \texttt{src/collision\_events.py} forme une matrice creuse $R$
à trois entrées par événement, telle que $C=R^TR$. L'application
$y\mapsto R^T(Ry)$ évite toute matrice dense de taille modes par modes.
Aux très grandes énergies réduites, certains $d_\mu$ deviennent nuls
en arithmétique machine. Une action limite finie ne valide alors pas
l'inversibilité du changement de coordonnées.
La génération et le stockage des événements peuvent néanmoins devenir
coûteux : ce prototype ne résout pas à lui seul le problème de ressources.

Les facteurs sont évalués avec des différences d'énergies et des logarithmes
stables. Cela évite de perdre le coefficient de désintégration lorsque les
occupations thermiques deviennent très petites.
L'équation non linéaire est conservée séparément comme référence à
occupations modérées.

Les tests comparent la matrice à la dérivée de cette équation, y compris
avec une référence calculée à 100 chiffres puis comparée en double
précision avec des tolérances explicites. Cela ne signifie pas que
l'opérateur soit précis à 100 chiffres. Ils contrôlent énergie, impulsion
normale, dissipation Umklapp, parité, duplication et limites numériques.
Le flux se simplifie en
\[
 F(n)=\Gamma[n_i(1+n_j+n_k)-n_jn_k].
\]
Il est quadratique : une différence centrée n'a pas d'erreur de troncature
en arithmétique exacte. Les erreurs observées lors de ce test sont donc
des erreurs d'arrondi ou de formulation, plutôt qu'une convergence
quadratique à rechercher systématiquement.

Les indices répétés sont exclus de cette API. Leur vecteur
stoechiométrique contient une multiplicité deux ; les facteurs
microscopiques supplémentaires doivent être vérifiés avant extension.
Les modes de fréquence exactement nulle et les poids irréductibles
inégaux sont également exclus.

\section{Passage aux données réelles}
Une famille complète d'entrées publiques AlN a été téléchargée sur le
SSD D, depuis une révision fixée du dépôt Phonon Olympics :
structure, correction polaire, constantes harmoniques et cubiques.
Les tailles et empreintes sont enregistrées. Ces fichiers ne sont
ni des événements déjà validés ni une trajectoire en température.

Le prochain import doit documenter : cartes de modes et d'inversion,
unités du couplage cubique, poids et multiplicités, canaux, vecteurs
réciproques, traitement de la conservation d'énergie et normalisation
temporelle. Un petit export des interactions est préférable comme
premier contrôle à une utilisation directe de la matrice de conductivité.
Aucun calcul de transport de l'AlN ni aucune validation expérimentale
n'est annoncé dans cette note.

\section{Sources et portée}
Fugallo et al., \emph{Ab initio variational approach for evaluating lattice
thermal conductivity}, arXiv:1212.0470v2, équations (4), (13), annexe A :
structure signée et conventions. Leur discussion de l'élargissement
souligne que le bilan détaillé devient approximatif.
Les références exactes, les accès et les réserves sont conservés dans
le rapport bibliographique de la campagne, publié avec les quatre branches,
les quatre audits et les expériences sous
\texttt{theory/aln/events/}.

La décomposition par événements est un cadre connu. Cette note fournit
une validation limitée et reproductible, sans revendication de nouveauté.
La principale raison pour laquelle une extension pourrait échouer est
une différence de comptage, de coordonnées ou d'intégration énergétique
entre ces événements idéaux et l'export matériel.
\end{document}
